\nonstopmode{}
\documentclass[letterpaper]{book}
\usepackage[times,inconsolata,hyper]{Rd}
\usepackage{makeidx}
\makeatletter\@ifl@t@r\fmtversion{2018/04/01}{}{\usepackage[utf8]{inputenc}}\makeatother
% \usepackage{graphicx} % @USE GRAPHICX@
\makeindex{}
\begin{document}
\chapter*{}
\begin{center}
{\textbf{\huge Package `gpcp'}}
\par\bigskip{\large \today}
\end{center}
\ifthenelse{\boolean{Rd@use@hyper}}{\hypersetup{pdftitle = {gpcp: Genomic Prediction of Cross Performance}}}{}
\ifthenelse{\boolean{Rd@use@hyper}}{\hypersetup{pdfauthor = {Marlee Labroo; Christine Nyaga; Lukas Mueller}}}{}
\begin{description}
\raggedright{}
\item[Type]\AsIs{Package}
\item[Title]\AsIs{Genomic Prediction of Cross Performance}
\item[Version]\AsIs{0.1.1}
\item[Author]\AsIs{Marlee Labroo [aut], Christine Nyaga [cre, aut], Lukas Mueller [aut]}
\item[Maintainer]\AsIs{Christine Nyaga }\email{cmn92@cornell.edu}\AsIs{}
\item[Description]\AsIs{This function performs genomic prediction of cross performance…}
\item[License]\AsIs{GPL (>= 3)}
\item[Encoding]\AsIs{UTF-8}
\item[LazyData]\AsIs{true}
\item[LinkingTo]\AsIs{Rcpp, RcppArmadillo}
\item[Imports]\AsIs{Rcpp, dplyr, sommer (>= 4.0.2), AGHmatrix, Matrix,
VariantAnnotation, tools, magrittr, methods}
\item[Suggests]\AsIs{knitr, rmarkdown, testthat (>= 3.0.0), snpStats, BiocManager}
\item[RoxygenNote]\AsIs{7.3.2}
\item[VignetteBuilder]\AsIs{knitr}
\item[Depends]\AsIs{R (>= 2.10)}
\item[Config/testthat/edition]\AsIs{3}
\end{description}
\Rdcontents{Contents}
\HeaderA{phenotypeFile}{Example Phenotype Data}{phenotypeFile}
\keyword{datasets}{phenotypeFile}
%
\begin{Description}
This is a sample phenotype dataset used for genomic prediction.
\end{Description}
%
\begin{Usage}
\begin{verbatim}
phenotypeFile
\end{verbatim}
\end{Usage}
%
\begin{Format}
A data frame with 24 columns:
\begin{description}

\item[ATW] Description of ATW
\item[AUDPC\_YAD] Area Under Disease Progress Curve for YAD
\item[AUDPC\_YMV] Area Under Disease Progress Curve for YMV
\item[Accession] Genotype IDs for each individual
\item[Block] Block information
\item[DMC] Dry Matter Content values
\item[Design] Experimental design
\item[LOC] Location of the trials
\item[NPH] Number of Plants Harvested
\item[OXBI] Oxidation Index
\item[Oxint180Minutes] Oxidation intensity after 180 minutes
\item[PLOT] Plot number
\item[REP] Replication number
\item[Settweight] Weight of the planting setts
\item[TTNPL] Total Tuber Number per Plant
\item[TTWPL] Total Tuber Weight per Plant
\item[Trial] Trial name or ID
\item[Vigor] Plant vigor score
\item[YIELD] Yield values
\item[Year] Year of the experiment
\item[Yield.per.plot..kg.] Yield per plot in kilograms
\item[Yield\_udj] Unadjusted Yield
\item[rAUDPC\_YAD] Relative AUDPC for YAD
\item[rAUDPC\_YMV] Relative AUDPC for YMV

\end{description}

\end{Format}
%
\begin{Source}
Generated for the gpcp package example
\end{Source}
%
\begin{Examples}
\begin{ExampleCode}
data(phenotypeFile)
head(phenotypeFile)
\end{ExampleCode}
\end{Examples}
\HeaderA{runGPCP}{Genomic Prediction of Cross Performance This function performs genomic prediction of cross performance using genotype and phenotype data.}{runGPCP}
%
\begin{Description}
Genomic Prediction of Cross Performance
This function performs genomic prediction of cross performance using genotype and phenotype data.
\end{Description}
%
\begin{Usage}
\begin{verbatim}
runGPCP(
  phenotypeFile,
  genotypeFile = NA,
  genotypeData = NA,
  genotypes,
  traits,
  weights = NA,
  userSexes = "",
  userFixed = NA,
  userRandom = NA,
  Ploidy = NA,
  NCrosses = NA
)
\end{verbatim}
\end{Usage}
%
\begin{Arguments}
\begin{ldescription}
\item[\code{phenotypeFile}] A data frame containing phenotypic data, typically read from a CSV file.

\item[\code{genotypeFile}] Path to the genotypic data, either in VCF or HapMap format.

\item[\code{genotypeData}] A dataframe containing genotypic data if genotypeFile not provided

\item[\code{genotypes}] A character string representing the column name in the phenotype file for the genotype IDs.

\item[\code{traits}] A string of comma-separated trait names from the phenotype file.

\item[\code{weights}] A numeric vector specifying weights for the traits.

\item[\code{userSexes}] A string representing the column name corresponding to the individuals' sexes.

\item[\code{userFixed}] A string of comma-separated fixed effect variables.

\item[\code{userRandom}] A string of comma-separated random effect variables.

\item[\code{Ploidy}] An integer representing the ploidy level of the organism.

\item[\code{NCrosses}] An integer specifying the number of top crosses to output.
\end{ldescription}
\end{Arguments}
%
\begin{Value}
A data frame containing predicted cross performance.
\end{Value}
%
\begin{Examples}
\begin{ExampleCode}
# Load phenotype data from CSV
# Diploid pipeline
phenotypeFile <- read.csv(system.file("extdata", "phenotypeFile.csv", package = "gpcp"))
genotypeFile <- system.file("extdata", "genotypeFile_Chr9and11.vcf", package = "gpcp")
finalcrosses <- runGPCP(
    phenotypeFile = phenotypeFile,
    genotypeFile = genotypeFile,
    genotypes = "Accession",
    traits = "YIELD,DMC",
    weights = c(3, 1),
    userFixed = "LOC,REP",
    Ploidy = 2,
    NCrosses = 150
)
message(finalcrosses)
 #PolyPLoid Pipeline
 # 1) load example data from the package
data(DT_polyploid, package = "sommer")
DT <- DT_polyploid
GT <- GT_polyploid
MP <- MP_polyploid
# 2) convert A/T/C/G strings to numeric codes
numo <- sommer::atcg1234(data = GT, ploidy = 4)

# 3) find the set of individuals common to genotypes and phenotypes
common <- intersect(DT$Name, rownames(numo$M))
marks  <- numo$M[common, , drop = FALSE]
pheno2 <- as.data.frame(DT[match(common, DT$Name), ])
# 4) call runGPCP with ploidy = 4
result4x <- suppressWarnings(
  runGPCP(
    phenotypeFile = pheno2,
    genotypeData   = marks,
    genotypes      = "Name",
    traits         = c("total_yield", "tuber_length"),
    weights        = c(3, 1),
    Ploidy         = 4,
    NCrosses       = 100
  )
)
\end{ExampleCode}
\end{Examples}
\printindex{}
\end{document}
